\section*{अध्याय \ref{ch:b2:reaction-enthalpy} --- अभिक्रिया की एन्थैल्पी}

\begin{solution}{exo:b2:reaction-enthalpy:1}
$\Delta_r H^\circ = 3(-393.51) + 4(-285.83) - (-104.7) - 5(0) =
\qty{-2219.2}{kJ/mol}$, ऋणात्मक: \omterm{def:b2:reaction-enthalpy:exothermic}{ऊष्माक्षेपी}। (मापी गई दहन
एन्थैल्पी भी यही है, \qty{-2219.2}{kJ/mol}।)
\end{solution}

\begin{solution}{exo:b2:reaction-enthalpy:2}
मेथेन: $890.3/16.04 = \qty{55.5}{kJ/g}$, अतः प्रति किलोग्राम \qty{55.5}{MJ}।
डाइहाइड्रोजन: $285.83/2.016 = \qty{141.8}{kJ/g}$,
\qty{141.8}{MJ/kg}: प्रति किलोग्राम लगभग 2.6 गुना अधिक (परंतु प्रति
लीटर कहीं कम, क्योंकि गैस हल्की है)।
\end{solution}

\begin{solution}{exo:b2:reaction-enthalpy:3}
$\Delta\nu_{\text{गैस}} = 6 - 7.5 = -1.5$, अतः $\Delta_r H = \Delta_r U +
\Delta\nu_{\text{गैस}}RT = -3264 - 1.5 \times 8.314 \times 298.15 \times
10^{-3} = \qty{-3267.7}{kJ/mol}$, जो संभवन एन्थैल्पियों से परिकलित मान,
\qty{-3267.6}{kJ/mol}, से मेल खाता है।
\end{solution}

\begin{solution}{exo:b2:reaction-enthalpy:4}
लक्ष्य पहली अभिक्रिया ऋण दूसरी का दुगुना है:
$\Delta_r H^\circ = -393.5 - 2(-283.0) = \qty{+172.5}{kJ/mol}$। यह
\omterm{def:b2:reaction-enthalpy:exothermic}{ऊष्माशोषी} है: गर्म कोक कार्बन डाइऑक्साइड को कार्बन मोनॉक्साइड में
बदलते हुए ठंडा होता है।
\end{solution}

\begin{solution}{exo:b2:reaction-enthalpy:5}
टूटे आबंध: C--H और Cl--Cl, $416 + 243 = \qty{659}{kJ/mol}$; बने
आबंध: C--Cl और H--Cl, $350 + 432 = \qty{782}{kJ/mol}$: आकलन
\qty{-123}{kJ/mol}। संभवन एन्थैल्पियों से: $-83.7 - 92.3 + 74.9 =
\qty{-101.1}{kJ/mol}$। \qty{22}{kJ/mol} का अंतर माध्य C--Cl
मान से आता है, जो कई अणुओं का औसत है, और मेथेन में टूटे C--H आबंध से,
जो चारों के माध्य से प्रबल है।
\end{solution}

\begin{solution}{exo:b2:reaction-enthalpy:6}
\qty{298}{K} पर: $-241.83 + 285.83 = \qty{44.00}{kJ/mol}$। किरखॉफ़ से,
$\Delta C_p = 33.59 - 75.35 = \qty{-41.76}{J/(K.mol)}$ के साथ:
\qty{400}{K} पर $44.00 - 0.04176 \times 101.85 = \qty{39.75}{kJ/mol}$। सारणियाँ
$-242.85 + 282.59 = \qty{39.74}{kJ/mol}$ देती हैं: इस अंतराल में स्थिर ऊष्मा
धारिताएँ उत्कृष्ट हैं।
\end{solution}

\begin{solution}{exo:b2:reaction-enthalpy:7}
आयनन: $4.341 \times 96.485 = \qty{418.8}{kJ/mol}$; संलग्नन:
$-3.613 \times 96.485 = \qty{-348.6}{kJ/mol}$। \omterm{def:b2:reaction-enthalpy:lattice-enthalpy}{जालक एन्थैल्पी} $= 436.7 +
89.0 + 121.3 + 418.8 - 348.6 = \qty{717}{kJ/mol}$, जो सोडियम क्लोराइड के
\qty{787}{kJ/mol} से कम है: \ce{K+} आयन \ce{Na+} से बड़ा है,
आयन एक-दूसरे से अधिक दूर हैं और कम आकर्षित करते हैं।
\end{solution}

\begin{solution}{exo:b2:reaction-enthalpy:8}
कैलोरीमीटर की ऊष्मा धारिता: $C = 2.00/0.418 = \qty{4.785}{kJ/K}$।
मुक्त ऊष्मा: $4.785 \times 0.583 = \qty{2.789}{kJ}$, $\xi =
\qty{0.0500}{mol}$ के लिए (अम्ल सीमांत है, क्षारक थोड़ा आधिक्य में):
$\Delta_r H = -2.789/0.0500 = \qty{-55.8}{kJ/mol}$।
\end{solution}

\begin{solution}{exo:b2:reaction-enthalpy:9}
\ce{C8H18(l) + 25/2O2(g) -> 8CO2(g) + 9H2O(l)}: $\Delta_r H^\circ =
8(-393.51) + 9(-285.83) + 250.3 = \qty{-5470.3}{kJ/mol}$; $M =
\qty{114.2}{g/mol}$ के साथ, \qty{47.9}{MJ/kg}। कार्बन डाइऑक्साइड: $8 \times 44.01 =
\qty{352.1}{g}$ प्रति \qty{5.470}{MJ}, अर्थात् \qty{64.4}{g/MJ}, जबकि
मेथेन के लिए $44.01/0.8903 = \qty{49.4}{g/MJ}$: मेथेन में प्रति कार्बन अधिक
हाइड्रोजन है, और हाइड्रोजन का दहन कार्बन डाइऑक्साइड के बिना ऊष्मा मुक्त करता है।
\end{solution}

\begin{solution}{exo:b2:reaction-enthalpy:10}
\ce{CH4 + 2O2 -> CO2 + 2H2O(g)} के उत्पाद: $C_p = 37.13 + 2 \times 33.59 =
\qty{104.3}{J/K}$; $q = \qty{802.3}{kJ}$; $T_f = 298 + 802\,300/104.3 \approx
\qty{7990}{K}$, वायु वाले मान का लगभग तीन गुना, क्योंकि कोई नाइट्रोजन
ऊष्मा में अपना हिस्सा नहीं लेती। ऑक्सीजन की वास्तविक ज्वाला कहीं ठंडी होती है: \qty{7990}{K} से बहुत पहले कार्बन डाइऑक्साइड और जल
कार्बन मोनॉक्साइड, हाइड्रोजन, ऑक्सीजन और परमाणुओं में वियोजित हो जाते हैं, ये \omterm{def:b2:reaction-enthalpy:exothermic}{ऊष्माशोषी}
अभिक्रियाएँ ऊष्मा का अधिकांश भाग अवशोषित कर लेती हैं; और ऊष्मा धारिताएँ
$T$ के साथ बढ़ती हैं।
\end{solution}

\begin{solution}{exo:b2:reaction-enthalpy:11}
$\Delta_r H^\circ(700) = \Delta_r H^\circ(298) + \int_{298}^{700}(\Delta_r a +
\Delta_r b\,T)\,\dd T = -91.88 + [-60.5 \times 401.85 + 0.0270 \times
(700^2 - 298.15^2)] \times 10^{-3} = -91.88 - 24.31 + 10.83 =
\qty{-105.4}{kJ/mol}$, जो सारणियों (\qty{-105.2}{kJ/mol}) से
\qty{0.2}{kJ/mol} के भीतर है, जबकि स्थिर-$C_p$ आकलन
\qty{4.5}{kJ/mol} हटकर था।
\end{solution}

\begin{solution}{exo:b2:reaction-enthalpy:12}
$\Delta_r H^\circ = 2 \times 90.29 = \qty{+180.6}{kJ/mol}$ (\omterm{def:b2:reaction-enthalpy:exothermic}{ऊष्माशोषी})।
$\Delta_r C_p^\circ = 2(29.85) - 29.12 - 29.38 = \qty{1.2}{J/(K.mol)}$, अतः
\qty{298}{K} और \qty{2000}{K} के बीच \omterm{def:b2:reaction-enthalpy:reaction-quantity}{अभिक्रिया एन्थैल्पी} लगभग
$1.2 \times 1702 = \qty{2}{kJ/mol}$ बदलती है, एक प्रतिशत। \omterm{def:b2:reaction-enthalpy:exothermic}{ऊष्माशोषी}
अभिक्रिया उच्च ताप पर अनुकूलित होती है (\cref{ch:b2:equilibrium-shifts}):
ठंडी वायु लगभग कोई नाइट्रोजन मोनॉक्साइड नहीं बनाती, सबसे गर्म ज्वालाएँ
सबसे अधिक बनाती हैं, इसीलिए दहन का ताप इस प्रदूषक के विरुद्ध
एक उपाय है।
\end{solution}

\begin{solution}{pb:b2:reaction-enthalpy:1}
\textbf{1.} \ce{C4H10(g) + 13/2O2(g) -> 4CO2(g) + 5H2O(l)}।
\textbf{2.} $\Delta_r H^\circ = 4(-393.51) + 5(-285.83) + 125.6 =
\qty{-2877.6}{kJ/mol}$।
\textbf{3.} जल वाष्प के साथ, $4(-393.51) + 5(-241.83) + 125.6 =
\qty{-2657.6}{kJ/mol}$; अंतर, \qty{220.0}{kJ}, पाँच मोल जल के वाष्पन की
एन्थैल्पी है (हर मोल के लिए \qty{44.0}{kJ/mol})।
\textbf{4.} $2877.6/58.12 = \qty{49.5}{kJ/g}$ और $2657.6/58.12 =
\qty{45.7}{kJ/g}$।
\textbf{5.} दूसरा: जल वाष्प के रूप में निकल जाता है और उसकी संघनन ऊष्मा
नष्ट हो जाती है।
\textbf{6.} $230 \times 45.7 = \qty{1.05e4}{kJ}$, लगभग \qty{10.5}{MJ}।
\textbf{7.} नियत आयतन पर प्राप्त ऊष्मा $\Delta U$ है (दाब-कार्य
नहीं होता): बम $\Delta_r U$ मापता है।
\textbf{8.} $\Delta\nu_{\text{गैस}} = 4 - 1 - 6.5 = -3.5$ (जल
द्रव है)।
\textbf{9.} $\Delta_r U = \Delta_r H - \Delta\nu_{\text{गैस}}RT = -2877.6 +
3.5 \times 2.479 = \qty{-2868.9}{kJ/mol}$।
\textbf{10.} $\xi = 0.500/58.12 = \qty{8.60e-3}{mol}$; ऊष्मा $8.60 \times
10^{-3} \times 2868.9 = \qty{24.68}{kJ}$; $\Delta T = 24.68/10.00 =
\qty{2.468}{K}$।
\textbf{11.} ऊष्मा $10.00 \times 2.461 = \qty{24.61}{kJ}$: $\Delta_r U =
-24.61/(8.603 \times 10^{-3}) = \qty{-2860.6}{kJ/mol}$ और $\Delta_r H =
-2860.6 - 8.7 = \qty{-2869.3}{kJ/mol}$, निरपेक्ष मान में सारणियों से
\qty{0.29}{\percent} कम: कैलोरीमीटर से ऊष्मा का क्षय या थोड़ा
अपूर्ण दहन।
\textbf{12.} $n = 1000/18.02 = \qty{55.5}{mol}$; $Q = 55.5 \times 75.35
\times 80 = \qty{334.5}{kJ}$।
\textbf{13.} मुक्त ऊष्मा $16.0 \times 45.72 = \qty{731.5}{kJ}$; दक्षता
$334.5/731.5 = 0.46$।
\textbf{14.} बर्तन के चारों ओर से निकलने वाली गर्म दहन गैसों में,
ज्वाला द्वारा गर्म की गई वायु और बर्तन की दीवारों में, और विकिरण में।
\textbf{15.} $230 \times 45.72 \times 0.457/334.5 = \qty{14.4}{L}$।
\textbf{16.} \ce{O2} के \qty{6.5}{mol} के साथ \ce{N2} के $6.5 \times 78.08/20.95 =
\qty{24.23}{mol}$ और \ce{Ar} के $6.5 \times 0.93/20.95 = \qty{0.289}{mol}$
आते हैं।
\textbf{17.} \qty{4}{mol} \ce{CO2}, \qty{5}{mol} \ce{H2O(g)},
\qty{24.23}{mol} \ce{N2}, \qty{0.289}{mol} \ce{Ar}।
\textbf{18.} ज्वाला रुद्धोष्म और समदाबी है, अतः $\Delta H = 0$। ऐसे
पथ के अनुदिश जो \qty{298}{K} पर अभिक्रिया ($\Delta H_1 =
\qty{-2657.6}{kJ}$) और उसके बाद इन उत्पादों के तापन से बना है, $\Delta H_2 =
-\Delta H_1 = \qty{+2657.6}{kJ}$: पूरी ऊष्मा गैसों में जाती है।
\textbf{19.} $C_p = 4(37.13) + 5(33.59) + 24.23(29.12) + 0.289(20.79) =
\qty{1028}{J/K}$।
\textbf{20.} $T_f = 298 + 2\,657\,600/1028 = \qty{2883}{K}$।
\textbf{21.} $T_f = 2300 + 100 \times (2657.6 - 2515.7)/(2655.7 - 2515.7) =
\qty{2401}{K}$।
\textbf{22.} गैसों की ऊष्मा धारिताएँ ताप के साथ बढ़ती हैं (\qty{298}{K} और
\qty{2400}{K} के बीच लगभग एक-चौथाई): कमरे के ताप वाले
मान गैसों का तापन बहुत सस्ता आँकते हैं।
\textbf{23.} \qty{2000}{K} से ऊपर कार्बन डाइऑक्साइड और जल आंशिक रूप से वियोजित होते हैं
(\omterm{def:b2:reaction-enthalpy:exothermic}{ऊष्माशोषी}), ज्वाला विकिरण करती है, और वायु के साथ मिश्रण कभी पूर्ण नहीं होता।
\textbf{24.} वायु में ब्यूटेन का \omterm{def:b2:reaction-enthalpy:flame-temperature}{रुद्धोष्म ज्वाला ताप}
$\boldsymbol{T_{\text{ad}} \approx \qty{2.40e3}{K}}$ है।
\end{solution}
